% --------------------------------------------------------------------
% --------------------------------------------------------------------
% Modelo de Trabalho Acadêmico utilizando classe repUERJ para elaboração
% de teses, dissertação e trabalhos monográficos em geral.
%
% Este arquivo está editado na codificação de caracteres UTF-8.
%
% As referencia estão baseadas no modelo bibtex e citação em autor-data
%
% Este modelo foi criado por Dr. Luís Fernando de Oliveira.
% Professor Adjunto do Departamento de Física Aplicada e Termodinâmica
% Instituto de Física Armando Dias Tavares
% Universidade do Estado do Rio de Janeiro - UERJ
%
% A classe repUERJ.cls foi criada a partir do código original 
% disponibilizado pelo grupo CódigoLivre (coordenado por
% Gerald Weber).
% Foram feitas adequações para implementação das normas de elaboração
% de teses e dissertações da UERJ.
%
% Os estilos repUERJformat.sty codificam os elementos pré-textuais e
% pós-textuais.
% O estilo repUERJpseudocode.sty codifica a elaboração de algoritmos
% utilizando um glossário desenvolvido por mim (Luís Fernando), o mesmo
% usado em meu curso de Física Computacional.
%
% Todo este material está disponível também no meu site
%      http://sites.google.com/site/deoliveiralf
%
% As normas da UERJ para elaboração de teses e dissertações pode ser
% obtidas no documento disponível no site
%      http://www.bdtd.uerj.br/roteiro_uerj_web.pdf
%
% Agradecimentos ao grupo da Rede Sirius - BDTD e à Biblioteca Setorial
% da Física.
% --------------------------------------------------------------------
% --------------------------------------------------------------------
%
\documentclass[a4paper,12pt,oneside,onecolumn,final,fleqn]{config/repUERJ}
% ---
% Pacotes fundamentais 
% ---
\usepackage[brazil,british]{babel}  % adequação para o português Brasil
\usepackage[utf8]{inputenc} % Determina a codificação utilizada
                            % (conversão automática dos acentos)
\usepackage{makeidx}        % Cria o índice
\usepackage{hyperref}       % Controla a formação do índice
\usepackage{indentfirst}    % Endenta o primeiro paragrafo de
                            % cada seção.
\usepackage{graphicx}       % Inclusão de gráficos
\graphicspath{images}
\usepackage{subfig}
\usepackage{amsmath}        % pacote matemático
\usepackage{amssymb}        % pacote extra
\usepackage{mathrsfs}       % pacote extra
\usepackage{tikz}	    % awesome drawings
\usetikzlibrary{calc}               % to get node width and height
\usetikzlibrary{positioning}        % to use the relative positioning commands
\usetikzlibrary{arrows.meta}
\usetikzlibrary{mindmap}
\usetikzlibrary{shapes.geometric}   % to use diamond shape
\usetikzlibrary{backgrounds}

% ---
% Pacote auxiliar para as normas da UERJ
% ---
\usepackage[frame=no,algline=yes,font=default]{config/repUERJformat}
\usepackage{config/repUERJpseudocode}
% ---
% Pacotes de citacoes
% ---
\usepackage[alf]{abntex2cite}

% ********************************************************************
% ********************************************************************
% Informações de autoria e institucionais
% ********************************************************************
% ********************************************************************

%---------------------------------------------------------------------
% Imagens pretextuais (precisam estar no mesmo diretório deste arquivo .tex)
%---------------------------------------------------------------------

\logo{images/logo_uerj_cinza.png}
\marcadagua{images/marcadagua_uerj_cinza.png}{1}{160}{255}

%---------------------------------------------------------------------
% Informações da instituição
%---------------------------------------------------------------------

\instituicao{Universidade do Estado do Rio de Janeiro}
            {Centro de Tecnologia e Ci\^{e}ncias} 
            {Instituto de Matem\'{a}tica e Estat\'{i}stica} 
            {} 

%---------------------------------------------------------------------
% Informações da autoria do documento
%---------------------------------------------------------------------

\autor{Artur}
      {Chiaperini Grover}
      {A. C. G.}

      \titulo{Intelig\^{e}ncia Computacional Usando M\'{a}quinas de Boltzmann}
% se não for usar a quarta palavra chave, deixar o campo vazio: {}
\palavraschaves{m\'{a}quina de Boltzmann}
               {m\'{a}quina restrita de Boltzmann}
               {diverg\^{e}ncia de Kullback-Leibler}
               {entropia relativa}

\title{Computational Intelligence Using Boltzmann Machines}
\keywords{Boltzmann machine}
         {restricted Boltzmann machine}
         {Kullback-Leibler divergence}
         {relative entropy}

\orientador{Dra.}%
           {Roseli Suzi}{Wedemann}%
           {IME/CComp -- UERJ}

%coorientador é opcional
%\coorientador{Cargo Titulação} 
%             {Nome}{Sobrenome} 
%             {Unidade -- Instituição} 

%---------------------------------------------------------------------
% Grau pretendido (Doutor, Mestre, Bacharel, Licenciado) e Curso
%---------------------------------------------------------------------

\grau{Mestre}  
\curso{Curso}

% área de concentração é opcional
%\areadeconcentracao{área}

%---------------------------------------------------------------------
% Informações adicionais (local, data e paginas)
%---------------------------------------------------------------------

\local{Rio de Janeiro} 
\data{22}{Novembro}{2019} 

% ********************************************************************
% ********************************************************************
% Configurações de aparência do PDF final
% ********************************************************************
% ********************************************************************

% alterando o aspecto da cor azul
\definecolor{blue}{RGB}{41,5,195}
%\definecolor{apricot}{RGB}{251,206,177}

% informações do PDF
\hypersetup{%
  %backref=true,
  %pagebackref=true,
  %bookmarks=true,
  unicode=false,
  pdftitle={\UERJtitulo},
  pdfauthor={\UERJautor},
  pdfsubject={\UERJpreambulo},
  pdfkeywords={PALAVRAS}{CHAVES}{\chaveA}{\chaveB}{\chaveC}{\chaveD},
  pdfproducer={\packagename},       % producer of the document
  pdfcreator={\UERJautor},
  colorlinks=true,       % false: boxed links; true: colored links
  linkcolor=black,       % color of internal links blue
  citecolor=black,       % color of links to bibliography blue
  filecolor=black,       % color of file links magenta
  urlcolor=black,
  bookmarksdepth=4
}

% ********************************************************************
% ********************************************************************
% Início do documento
% ********************************************************************
% ********************************************************************
% ---
% compila o índice; se não for usar, comentar
% ---
\makeindex
% ---
% ********************************************************************
% ********************************************************************
\begin{document}
% ----------------------------------------------------------
% ELEMENTOS PRE-TEXTUAIS
% ----------------------------------------------------------
\frontmatter
% ----------------------------------------------------------
% Capa e a folha de rosto
% ----------------------------------------------------------
\capa%
\folhaderosto%
% ----------------------------------------------------------
% Inserir a ficha catalográfica
% ----------------------------------------------------------

% A biblioteca deverá providenciar a ficha catalográfica. Salve a ficha no formato PDF.
% Use o nome do arquivo PDF como argumento do comando. Exemplo: ficha catalográfica
% no arquivo 'ficha.pdf'

%\fichacatalografica{ficha.pdf}

% Enquanto não possuir a ficha catalográfica, use o comando sem argumentos...

\fichacatalografica{}

% ----------------------------------------------------------
% Folha de aprovação
% ----------------------------------------------------------
% membros da banca: máximo 6
\begin{folhadeaprovacao}
  \assinatura{Dra. Patr\'{i}cia Nunes da Silva}{IME/CComp -- UERJ}
  \assinatura{Dra. Zochil Gonz\'{a}lez Arenas}{IME/CComp -- UERJ}
%  \assinatura{terceiro membro titular da banca}{instituição}
% suplente só é incluído se efetivamente substitui um titular
%  \assinatura{primeiro membro suplente da banca}{instituição}
%  \assinatura{segundo membro suplente da banca}{instituição}
%  \assinatura{terceiro membro suplente da banca}{instituição instituição instituição instituição }
\end{folhadeaprovacao}
% ----------------------------------------------------------
% Dedicatória
% ----------------------------------------------------------
%\pretextualchapter{Dedicatória}
%
%\vfill
%Texto da dedicatória
% ----------------------------------------------------------
% Agradecimentos
% ----------------------------------------------------------
%\pretextualchapter{Agradecimentos}
%
%Texto de agradecimento
% ----------------------------------------------------------
% Epigrafe (opcional)
% ----------------------------------------------------------
\pretextualchapter{}

\vfill\
\begin{flushright}
    If no one died, it is just another story to be told.
    
    [Daniel Mirolhaum]
\end{flushright}
% ----------------------------------------------------------
% RESUMO
% ----------------------------------------------------------
\pretextualchapter{Resumo}

\referencia%

Texto do resumo em português.\\

\imprimirchaves%
% ----------------------------------------------------------
% Abstract
% ----------------------------------------------------------
\pretextualchapter{Abstract}

\reference%

Abstract in English.\\

\printkeys%
% ----------------------------------------------------------
% Listas de ilustrações e tabelas
% ----------------------------------------------------------
\listadefiguras%
%\listadetabelas%
% ----------------------------------------------------------
% Outras listas
% ----------------------------------------------------------
%\listadealgoritmos%
% ----------------------------------------------------------
% Lista de abreviaturas e siglas
% ----------------------------------------------------------
%\pretextualchapter{Lista de abreviaturas e siglas}
\pretextualchapter{List of abbreviations and acronyms}

\abreviatura{BM}{Boltzmann Machine}
\abreviatura{RBM}{Restricted Boltzmann Machine}
\abreviatura{SNN}{Stochastic Neural Network}
\abreviatura{ANN}{Artificial Neural Network}
\abreviatura{KL}{Kullback-Leibler}

%\abreviatura{sigla1}{por extenso}
%\abreviatura{sigla2}{por extenso}
%\abreviatura{sigla3}{por extenso}
% ----------------------------------------------------------
% Lista de simbolos
% ----------------------------------------------------------
%\pretextualchapter{Lista de símbolos}

%\simbolo{simbolo1}{significado e/ou valor}
%\simbolo{simbolo2}{significado e/ou valor}
%\simbolo{simbolo3}{significado e/ou valor}
% ----------------------------------------------------------
% Sumario
% ----------------------------------------------------------
\sumario%
% ----------------------------------------------------------
% ELEMENTOS TEXTUAIS
% ----------------------------------------------------------
\mainmatter%
%=====================================================================
\chapter*{Introduction}%
%=====================================================================
\label{ch:intro}%
Nowadays, modelling intelligent complex systems uses two main paradigms, commonly referred to as Symbolic and Connectionist, as basic guidelines for achieving your goals of creating intelligent machines and understanding human cognition.
These two approaches depart from different positions, each advocating advantages over the other in reproducing intelligent activity.
The traditional symbolic approach argues that the algorithmic manipulation of symbolic systems is an appropriate context for modelling cognitive processes.
On the other hand, connectionists restrict themselves to brain-inspired architectures and argue that this approach has the potential to overcome the rigidity of symbolic systems by more accurately modeling cognitive tasks that can only be solved, in the best case, approximately. In addition, connectionist paradigm focus on parallel models which have simple and uniform processing elements interconnected~\cite{bib:sun2001}.

Years of experimentation with both paradigms lead us to the conclusion that the solution lies between these two extremes, and that the approaches must be integrated and unified.
In order to establish a proper link between them, much remains to be researched.
If in the 1980s the discussion of intelligence was placed at the distinct poles of the symbolists and connectionists, today the connectionists are divided by the reductionist arguments of the structuralists.
For this structuralist current, the failure of the symbolists was due to the fact that their models despised brain architecture, and therefore connectionism must continue to explore more deeply the structural aspects of the thinking organ.

The connectionist and structuralist aspects are approached, respectively, through the paradigm of artificial neural networks and realistic models of the brain, within the area called Computational Neuroscience.
Through our models, we investigate ancient questions of Artificial Intelligence regarding the understanding of computability aspects of the human mind.

In this project we will continue with the study and implementation of artificial neural network known as Boltzmann Machine and its derivatives models, such as the Restricted Boltzmann Machine, which have been used to solve artificial intelligence problems, in areas such as: computer vision, encoder problem, statistical pattern recognition, speech recognition, and combinatorial optimization~\cite{bib:hertz1991}.
Boltzmann machines is a branch of stochastic methods, where randomness is present.~\cite{bib:duda2000}, it is also used to solve \textit{searching} and \textit{learning} problems~\cite{bib:hinton2010}.

% automatic speech (or voice) recognition, image recognition and treatment, natural language processing, bioinformatics, among many others.

%Our previous experience, both in the development of research in the field of artificial neural networks and general distributed processing and its technological applications, as well as in the pursuit of realistic models of brain biology, allows us to mature in the same direction of multidisciplinary research.

%Texto da introdução. Texto, texto texto \cite{bib:Amado1991}, texto \citeonline{bib:Amado1991}. Texto \citeauthoronline{bib:Andrade1997} em \citeyear{bib:Andrade1997}, texto \citeauthor{bib:Andrade1997},  texto.

%=====================================================================
\chapter{Boltzmann Machines}%
%=====================================================================
\label{ch:bm}%
EXPLICAR OS ELEMENTOS DE PROBABILIDADE COM OS QUAIS ESTAMOS LIDANDO:%
QUEM S\~{A}O AS VARI\'{A}VEIS ALEAT\'{O}RIAS E COMO VAMOS DOMIN\'{A}-LAS NO TEXTO,
COMENTAR QUE ESTAMOS TRATANDO COM VARI\'{A}VEIS DISCRETAS, COMO VAMOS IDENTIFICAR AS PROBABILIDADES,$\ldots$

EXPLICAR O QUE OS \'{I}NDICES REPResentam, e as vari\'{a}veis.

EXPLICAR OS TERMOS $\omega$ e $\phi$. $\omega$ \'{E} para os pesos, e $\phi$ para o bias de cada unidade.

EXPLICAR MINHA NOTA\c{C}\~{A}O sobre o sobre-escrito entre par\^{e}nteses!!!


In this chapter will expose the Boltzmann Machine theory. 
A few considerations regarding the notation used in this exposition is required before stepping forward into the main content. 
A single random variables is denoted by $\mathrm{x}$.
A vector of random variables of size $n$ is represented by $\mathrm{\mathbf{x}} = (\mathrm{x}_{1}, \mathrm{x}_{2}, \dots, \mathrm{x}_{n})$, where each $\mathrm{x}_{i}$, $i \in \{1, 2, \dots, n\}$, represents a single unit of the network. 
The value a random variable can assume is represented by $x$. 
Assuming a discrete scenario, the probability of a random variable $\mathrm{x}_{i}$ of assuming a certain value $x_{i}$ is $P(\mathrm{x}_{i} = x_{i})$. 


\section{Boltzmann Machines}

Boltzmann Machines (BM) are a type of stochastic neural networks (SNN) where the connections between units, which are described by $\omega$, are symmetrical, i.e., $\omega_{ij} = \omega_{ji}$ [HERTZ].
This kind of stochastic neural networks are capable of learning internal representation and to model an input distribution.
Boltzmann Machines were named after the Boltzmann distribution.
Due to its stochatics behaviour, the probability of the state of the system to be found in a certain configuration is given by previous mentioned distribution [HERTZ].
According to [MONTUFAR, 2018], BM can be seen as an extension of Hopfield networks to include hidden units.

Boltzamann Machines have visible and hidden units.
The visible units are linked to the external world and they correspond to the components of an observation. On the other hand, the hidden units do not have any connection outside of the network and model the dependencies between the components of the observations [FISCHER, 2012].
In BM, there is no connection restriction, this means that every unit, visible or hidden, can be connected to every other unit as in a complete graph, this pattern is not mandatory as some of the connections may not exists depending on the network layout.
\begin{figure}[h!tbp]{\textwidth}
  \centering
  \caption{Boltzmann machine diagram.}%
  \label{fig:bm-diagram}%
  \tikzstyle{visible} = [draw, circle, fill=white]
  \tikzstyle{hidden} = [draw, circle, fill=gray!30]
  \begin{tikzpicture}
    \node[style=visible] (a) at (0, 2.5) {$\mathrm{x}_{1}$};
    % \node[style=visible] (b) at +(0: 1.5) {$\mathrm{x}_{2}$};
    % \node[style=hidden] (c) at + (60: 1.5) {$\mathrm{x}_{3}$};
    \node[style=visible] (b) at (-2.16, 1.25) {$\mathrm{x}_{2}$};
    \node[style=visible] (c) at (-2.16, -1.25) {$\mathrm{x}_{3}$};
    \node[style=visible] (d) at (0, -2.5) {$\mathrm{x}_{4}$};
    \node[style=hidden]  (e) at (2.16, -1.25) {$\mathrm{x}_{5}$};
    \node[style=hidden]  (f) at (2.16, 1.25) {$\mathrm{x}_{6}$};
    \foreach \from/\to in {a/b, a/c, a/d, a/e, a/f, b/c, b/d, b/e, b/f, c/d, c/e, c/f, d/e, d/f, e/f}
      \draw [thick] (\from) -- (\to);   
  \end{tikzpicture}
  \legend{Gray circles represent the hidden units of the Boltzmann Machine, while the white circles are the visible units.}
  \source{Author.}
\end{figure}


Training Boltzmann Machines means finding the right connection between the units.

Boltzmann Machines (BM) are stochastics neural networks with symmetric connections, i.e., $\omega_{ij} = \omega_{ji}$.
Boltzmann Machines use the Boltzmann distribution to determine the probability of the state of the system of the network.
BM ressambles the Hopfield networks with the inclusion of hidden units.
Finding the right connections between the hidden units without knowing it from the training patterns what the hidden units represent is part of the solving the Boltzmann Machine problem.

Units $x_{i}$ in BM are split into two kinds: visible and hidden units.
The visible units have connection to the outside world and are the units that receive the data input.
On the other hand, the hidden units do not have any connection to the outside of the network and they are resposible to find the data relation from the input.
In a BM, the connections between units can be complete or not.
Regardless of how the connections are, every connection in a BM is symmetric.

BM are made of stochastics units $x_{i}$.
Stochastics units are random variables that can assume a binary value with a certain probability. 
We will consider that a random variable $\mathrm{x}_{i}$ can assume a value $x_{i} \in \{0, 1\}$, i. e.,
\begin{equation}
  x_{i} =
    \begin{cases}
      1 \; \text{with probability} \; g(h_{i}) \\
      0 \; \text{with probability} \; 1 - g(h_{i})
    \end{cases},
  \label{eq:stochastic-unit-values}
\end{equation}
where the probability is given by
\begin{equation}
  g(h_{i}) = \frac{1}{1 + e^{-2 \beta h_{i}}},
  \label{eq:stochastic-unit-prob}
\end{equation}
and
\begin{equation}
  h_{i} = \sum_{j}\omega_{ij}x_{j}.
  \label{eq:stochastic-unit-input}
\end{equation}

Due to the symmetrical connections, there is an energy function give by
\begin{equation}
  H(\mathbf{x}) = - \sum_{i} \sum_{j} \omega_{ij}x_{i}x_{j} - \sum_{i} \phi_{i}x_{i},
  \label{eq:eq4}
\end{equation}
where $\mathbf{x} = (x_{1}, x_{2}, \dots, x_{n})$, and $n$ is equal to the number of units in the network, and the above energy function has minimum when there is a stable state characterised by
\begin{equation}
  x_{i} = sgn(h_{i}).
  \label{eq:eq5}
\end{equation}

The probability $P$ of finding the system in a given state $\mathbf{x}$ after the equilibrium is reached can be computed as follows:
\begin{equation}
  P(\mathbf{x}) = \frac{1}{Z} e^{-\beta H(\mathbf{x})},
  \label{eq:eq6}
\end{equation}
where
\begin{equation}
  Z = \sum_{\mathbf{x}'} e^{-\beta H(\mathbf{x}')}
  \label{eq:eq7}
\end{equation}
is the partition function.

The learning process of a Boltzmann Machine consists in ajusting the connections $\omega_{ij}$ in such a way that the state of the visible units have a particular desired probability distribution.

Let us identify the state of the visible units by an index $v$ and the state of the hidden units by an index $h$.
Considering a system which has N visible units and K hidden units, the whole system have $2^{N + K}$ possibilities of states in which it can be found.

The joint probability $P_{vh}$ is the probability of finding the visible and hidden units in the states $v$ and $h$, respectively.
This probabiblity measument is given by the Boltzmann distribution:
\begin{equation}
  P_{vh} = \frac{e^{-\beta H_{vh}}}{Z},
  \label{eq:eq8}
\end{equation}
where
\begin{equation}
  Z = \sum_{u} \sum_{k} e^{-\beta H_{uk}},
  \label{eq:eq9}
\end{equation}
and
\begin{equation}
  H_{vh} = - \sum_{i} \sum_{j} \omega_{ij} x^{(vh)}_{i} x^{(vh)}_{j} - \sum_{i} \phi_{i} x^{(vh)}_{i}.
  \label{eq:eq10}
\end{equation}

As metioned above, the problem a Boltzmann Machine is trying to solve is determining the connections $\omega_{ij}$ between units such that the visible units have a certain probability distribution.
In order to do that, we need to find the marginal probability of the state $v$ in which the visible units are found regardless of the state $h$ of the hidden units. The marginal probability $P_{v}$ is given by
\begin{equation}
  P_{v} = \sum_{h} P_{vh} = \sum_{h} \frac{e^{-\beta H_{vh}}}{Z}.
  \label{eq:marginal_prob}
\end{equation}

Although we know that $P_{v}$ is a function of the connections $\omega_{ij}$, and that this is the probability of finding the visible units in the state $v$. We want the states to have a certain probability $Q_{v}$, i.e., a desired probability.
This means that ideally we would like to match the empirical distribution of the data, even though we do not have access to the correct distribution, only to what the observed data has given us as an input to training the model.

One way to evaluate the difference between two probability distribution, for example, $P_{v}$ and $Q_{v}$, is using the Kullback-Leibler divergence, which can also be referred to relative entropy, $E$, which will be our cost function.
(EXPLICAR FUN\c{C}\~{A}O DE CUSTO e $D_{KL}$!!!).
\begin{equation}
  E = \sum_{v} Q_{v} \ln{\left(\frac{Q_{v}}{P_{v}}\right)}.
  \label{eq:DKL}
\end{equation}

The relative entropy $E$ has the property of always being equal or greater than zero.
It reaches zero only if $P_{v} = Q_{v}$, which means that we are able to retrieve the exactly probability distribution of the input data at the visible units.
\begin{equation}
  \begin{split}
    E & = \sum_{v} Q_{v} \ln{\left( \frac{Q_{v}}{P_{v}} \right)} \\
      & \geq \sum_{v} Q_{v} \left( 1 - \frac{P_{v}}{Q_{v}} \right) \\
      & = \sum_{v} \left( Q_{v} - P_{v} \right) \\
      & = \sum_{v} Q_{v} - \sum_{v} P_{v} = 1 - 1 \\
      & \Rightarrow E \geq 0.
  \end{split}
\end{equation}

From the gradient descent equation
\begin{equation}
  \Delta \omega_{ij} = -\eta \frac{\partial E}{\partial \omega_{ij}},
  \label{eq:gradient}
\end{equation}
where
\begin{equation}
  \begin{split}
    E & = \sum_{v} Q_{v} \ln{\left(\frac{Q_{v}}{P_{v}}\right)} \\
      & = \sum_{v} \left[ \ln{(Q_{v})} - \ln{(P_{v})} \right].
  \end{split}
  \label{eq:entropy1}
\end{equation}

In the following steps, we present the gradient descent derivation
\begin{equation}
  \begin{split}
    \Delta \omega_{ij} & = - \eta \frac{\partial E}{\partial \omega_{ij}} \\
                  & = - \eta \frac{\partial}{\partial \omega_{ij}} \left[ \sum_{v} Q_{v} \left( \ln{(Q_{v})} - \ln{(P_{v})} \right) \right] \\
                  & = \eta \frac{\partial}{\partial \omega_{ij}} \left[ \sum_{v} Q_{v} \ln{(P_{v})} \right] \\
                  & = \eta \sum_{v} Q_{v} \frac{\partial}{\partial \omega_{ij}} \left[ \ln{(P_{v})} \right] \\
                  \Rightarrow & \Delta \omega_{ij} = \eta \sum_{v} \frac{Q_{v}}{P_{v}} \frac{\partial P_{v}}{\partial \omega_{ij}}.
  \end{split}
  \label{eq:grad1}
\end{equation}

To continue with the computation of $\Delta \omega_{ij}$, we have to find the derivative of $\partial P_{v}/\partial \omega_{ij}$, from the marginal probability, equation~\ref{eq:marginal_prob},
\begin{equation}
  P_{v} = \frac{\sum_{h} e^{-\beta H_{vh}}}{\sum_{u} \sum_{k} e^{-\beta H_{u k}}},
  \label{eq:marginal_prob_expansion}
\end{equation}
thus the derivative of $P_{v}$ follows
\begin{equation}
  \begin{split}
    \frac{\partial P_{v}}{\partial \omega_{ij}} & = \frac{\partial}{\partial \omega_{ij}} \left[ \frac{\sum_{h} e^{-\beta H_{vh}}}{\sum_{u} \sum_{k} e^{-\beta H_{u k}}} \right] \\
    & = \frac{1}{\sum_{u} \sum_{k} e^{-\beta H_{u k}}} \sum_{h} (-\beta) e^{-\beta H_{vh}} \frac{\partial H_{vh}}{\partial \omega_{ij}} \\
    & - \sum_{h} e^{-\beta H_{vh}} \frac{1}{{\left( \sum_{u} \sum_{k} e^{-\beta H_{u k}} \right)}^{2}} \sum_{u} \sum_{k} e^{-\beta H_{u k}} (-\beta) \frac{\partial H_{u k}}{\partial \omega_{ij}}.
  \end{split}
  \label{eq:eq14}
\end{equation}

Following the equation~\ref{eq:eq14}, we need to compute the term $\partial H_{vh}/\partial \omega_{ij}$,
\begin{equation}
  \begin{split}
      \frac{\partial H_{vh}}{\partial \omega_{ij}} & = \frac{\partial}{\partial \omega_{ij}} \left[- \frac{1}{2} \sum_{m} \sum_{n} \omega_{mn} x^{(vh)}_{m} x^{(vh)}_{n} - \sum_{m} \phi_{mm} x^{(vh)}_{m} \right] \\
                                                   & = \frac{\partial}{\partial \omega_{ij}} \left[- \frac{1}{2} \sum_{m \neq i,j} \sum_{n \neq i,j} \omega_{mn} x^{(vh)}_{m} x^{(vh)}_{n} - \frac{1}{2} \omega_{ij} x^{(vh)}_{i} x^{(vh)}_{j} - \frac{1}{2} \omega_{ji} x^{(vh)}_{j} x^{(vh)}_{i} - \sum_{m} \phi_{m} x^{(vh)}_{m} \right],
  \end{split}
  \label{eq:eq15}
\end{equation}
as the connections between units are symmetric, i.e., $\omega_{ij} = \omega_{ji}$, then we can simplify equation~\ref{eq:eq15},
\begin{equation}
  \begin{split}
    \frac{\partial H_{vh}}{\partial \omega_{ij}} & = \frac{\partial}{\partial \omega_{ij}} \left[- \frac{1}{2} \sum_{m \neq i,j} \sum_{n \neq i,j} \omega_{mn} x^{(vh)}_{m} x^{(vh)}_{n} - \omega_{ij} x^{(vh)}_{i} x^{(vh)}_{j} - \sum_{m} \phi_{m} x^{(vh)}_{m} \right] \\
    & = \frac{\partial}{\partial \omega_{ij}} \left[- \frac{1}{2} \sum_{m \neq i,j} \sum_{n \neq i,j} \omega_{mn} x^{(vh)}_{m} x^{(vh)}_{n} \right] + \frac{\partial}{\partial \omega_{ij}} \left[- \omega_{ij} x^{(vh)}_{i} x^{(vh)}_{j} \right] + \frac{\partial}{\partial \omega_{ij}} \left[- \sum_{m} \phi_{m} x^{(vh)}_{m} \right] \\
    \Rightarrow & \frac{\partial H_{vh}}{\partial \omega_{ij}} = - x^{(vh)}_{i} x^{(vh)}_{j}.
  \end{split}
  \label{eq:eq16}
\end{equation}

Analagous to $\partial H_{vh}/\partial \omega_{ij}$, we have $H_{uk}$ derivative, which is
\begin{equation}
  \frac{\partial H_{uk}}{\partial \omega_{ij}} = - x^{(uk)}_{i} x^{(uk)}_{j}.
  \label{eq:eq17}
\end{equation}

Going back to equation~\ref{eq:eq14}, we can replace the derivatives of $H$, and solve the derivative of the marginal probability $P_{v}$,
\begin{equation}
  \begin{split}
    \frac{\partial P_{v}}{\partial \omega_{ij}} & = \frac{1}{Z} \sum_{h} e^{-\beta H_{vh}} (-\beta) (-x^{(vh)}_{i} x^{(vh)}_{j}) - \frac{1}{Z^{2}} \sum_{h} e^{-\beta H_{vh}} \sum_{u} \sum_{k} e^{-\beta H_{uk}} (-\beta) (-x^{(uk)}_{i} x^{(uk)}_{j}) \\
    & = \beta \left[ \sum_{h} \frac{e^{-\beta H_{vh}}}{Z} x^{(vh)}_{i} x^{(vh)}_{j} - \sum_{h} \frac{e^{-\beta H_{vh}}}{Z} \sum_{u} \sum_{k} \frac{e^{-\beta H_{uk}}}{Z} x^{(uk)}_{i} x^{(uk)}_{j}  \right] \\
    & = \beta \left[ \sum_{h} P_{vh} x^{(vh)}_{i} x^{(vh)}_{j} - P_{v} \sum_{u} \sum_{k} P_{uk} x^{(uk)}_{i} x^{(uk)}_{j} \right] \\
    \Rightarrow & \frac{\partial P_{v}}{\partial \omega_{ij}} = \beta \left[ \sum_{h} P_{vh} x^{(vh)}_{i} x^{(vh)}_{j} - P_{v} \langle x_{i} x_{j} \rangle  \right].
  \end{split}
  \label{eq:eq18}
\end{equation}

Given the derivative of $P_{v}$ in relation to $\omega_{ij}$, we can compute the learning term $\Delta \omega_{ij}$, from equation~\ref{eq:grad1},
\begin{equation}
  \begin{split}
    \Delta \omega_{ij} & = \eta \sum_{v} \frac{Q_{v}}{P_{v}} \beta \left[ \sum_{h} P_{vh} x^{(vh)}_{i} x^{(vh)}_{j} - P_{v} \langle x_{i} x_{j} \rangle \right]  \\
    & = \eta \beta \left[ \sum_{v} \sum_{h} \frac{Q_{v}}{P_{v}} P_{vh} x^{(vh)}_{i} x^{(vh)}_{j} - \sum_{v} \frac{Q_{v}}{P_{v}} P_{v} \langle x_{i} x_{j} \rangle \right] \\
    & = \eta \beta \left[ \sum_{v} \sum_{h} Q_{v} \frac{P_{vh}}{P_{v}} x^{(vh)}_{i} x^{(vh)}_{j} - \sum_{v} Q_{v} \langle x_{i} x_{j} \rangle \right] \\
    & = \eta \beta \left[ \sum_{v} \sum_{h} Q_{v} P_{h|v} x^{(vh)}_{i} x^{(vh)}_{j} - \langle x_{i} x_{j} \rangle \right].
  \end{split}
  \label{eq:learning_term2}
\end{equation}

In the above derivation, we have used the following relations,
\begin{equation}
  P_{h|v} = \frac{P_{vh}}{P_{v}},
  \label{eq:conditional_prob1}
\end{equation}
which is the conditional probability equation.
In our scenario this equation means that the probability distribution of the hidden units in state $h$ given the state $v$ of the visible units is the joint probability distribution of both states, $v$ and $h$, divided by the marginal probability distribution of the visible units in state $v$.

The second term in equation~\ref{eq:learning_term2}, is the average of units $i$ and $j$ over all combinations of states $v$ and $h$ of the system. 
In other words, we would have to compute all possible combination of states of visible and hidden units, $v$ and $h$, and then average over the specific units $i$ and $j$. 

The first term can be simplified by
\begin{equation}
  \sum_{v} Q_{v} \sum_{h} P_{h|v} x^{(vh)}_{i} x^{(vh)}_{j} = \sum_{v} Q_{v} \langle x_{i} x_{j} \rangle^{(v)} = \langle \langle x_{i} x_{j} \rangle^{(v)} \rangle.
  \label{eq:clamped_term}
\end{equation}

Then equation~\ref{eq:learning_term2} becomes
\begin{equation}
  \Delta \omega_{ij} = \eta \beta \left[ \langle \langle x_{i} x_{j} \rangle^{(v)} \rangle_{clamped} - \langle x_{i} x_{j} \rangle_{free} \right],
  \label{eq:learning_term3}
\end{equation}
it is important to notice that the subscripts \textit{clamped} means that we have to fix a certain $v$ state on the visible units otherwise the second term in the equation does not have a reference. 
On the other hand, the subscript \textit{free} identify the $\ldots$



%=====================================================================
%\chapter{T\'{i}tulo do cap\'{i}tulo 1}
%=====================================================================

%Texto do capítulo. Texto, texto, Figura~\ref{rotulo}. Texto Figura~\ref{outro.rotulo}\subref{subrotulo1}.

% \begin{figure ou table}[posição]{largura da figura}

%\begin{figure}[!ht]{6cm}
%  \caption{Título da figura.}\label{rotulo}
%  \includegraphics[width=\hsize]{images/logo_uerj_cor.jpg}
%  \legend{Texto da legenda.}
%  \source{Citação da fonte ou `O autor.'.}
%\end{figure}


%\begin{figure}[!ht]{11cm}
%  \caption{Título da figura.}\label{outro.rotulo}
%  \subfloat[][]{\label{subrotulo1}
%    \fbox{\includegraphics[width=0.45\hsize]{images/logo_uerj_cinza.png}}}\hfill
%  \subfloat[][]{\label{subrotulo2}
%    \fbox{\includegraphics[width=0.45\hsize]{images/marcadagua_uerj_cinza.png}}}\\
%  \subfloat[][]{\label{subrotulo3}
%    \fbox{\includegraphics[width=0.45\hsize]{images/logo_uerj_cor.jpg}}}\hfill
%  \legend{Texto da legenda. \subref{subrotulo1} Texto da imagem.
%          \subref{subrotulo2} Texto da imagem.
%          \subref{subrotulo3} Texto da imagem.}
%  \source{Citação da fonte ou `O autor'.}
%\end{figure}

%\begin{table}[!ht]{4cm}
%  \caption{Título da tabela.}\label{mais.rotulo}
%  \hfill\begin{tabular}{l|l}
%    \hline
%      X & Y\\
%    \hline
%      1,20 & 15,7\\
%      1,23 & 15,6\\
%      1,19 & 15,3\\
%      1,26 & 15,1\\
%      1,22 & 15,5\\
%      1,16 & 15,3\\
%      1,37 & 15,7\\
%    \hline
%  \end{tabular}\hfill
%  \legend{Texto da legenda.}
%  \source{Citação da fonte ou `O autor.'.}
%\end{table}

%=====================================================================
%\chapter{Título do capítulo 2}
%=====================================================================

%Texto do capítulo. Texto, texto Algoritmo~\ref{alg:modelo}. Texto.
%
%\begin{algorithm}[!ht]
%    \caption{Título do algoritmo.}\label{alg:modelo}
%    \begin{pseudocode}
%      \LinhaEmBranco%
%      \Documentacao%
%        \Titulo{Nome do algoritmo\\}
%        \Proposito{Propósito do algoritmo.\\}
%        \Metodo{Método utilizado no algoritmo.\\}
%        \Entradas{%
%          a, m: multiplicador e módulo\\
%          n0: semente\\
%          i: contador auxiliar\\
%        }
%        \Saidas{%
%          n: número aleatório\\
%        }
%        \Observacoes{Observações, restrições e requisitos.\\}
%        \Algoritmo{Identificação}
%          \Declarar{$a, m, i$}{numéricos}{}{}
%          \Declarar{$n0, n$}{numéricos}{}{}
%          \Ins{$m \leftarrow 13$}
%          \Ins{$n0 \leftarrow 1$}
%          \ParaDeAtePasso[para cada possível valor de `a']{$a$}{2}{$m-1$}{}
%            \Escrever{``a = '', $a$, ``: n = \{''}
%            \Ins[reinicia a geração com a semente n0]{$n \leftarrow n0$}
%            \ParaDeAtePasso{$i$}{0}{$m-1$}{}
%              \Ins[gerador de números aleatórios]{$n \leftarrow resto(a*n, m)$}
%              \SeEntao[se fim da sequencia $\dots$]{$n == n0$}
%                \Escrever{$n$,``\}''}
%                \Parar%
%              \Senao%
%                \Escrever{$n$}
%              \FimSe%
%            \FimPara%
%          \FimPara%
%        \Continua%
%    \end{pseudocode}
%\end{algorithm}
%
%\alglinenumbersoff%
%\begin{algorithm*}[!ht]
%    \caption{Título do algoritmo. (continuação)}
%    \begin{pseudocode*}
%        \LinhaEmBranco%
%        \Continuacao%
%%          \LinhaEmBranco
%          \Ins{$a \leftarrow 1$}
%          \Enquanto[comentário]{$a<10$}
%            \Escrever{$a$}
%            \Ins{$a \leftarrow a+1$}
%          \FimEnquanto%
%          \Ins{$a \leftarrow 1$}
%          \Repetir[comentário]
%            \Escrever{$a$}
%            \Ins{$a \leftarrow a+1$}
%          \AteQue{$a\ge10$}
%          \LinhaEmBranco%
%          \Ins{$a \leftarrow 1$}
%          \Fazer[comentário]
%            \Escrever{$a$}
%            \Ins{$a \leftarrow a+1$}
%          \Enquanto{$a<10$}
%        \FimAlgoritmo%
%      \FimDocumentacao%
%    \end{pseudocode*}
%\end{algorithm*}
%
%\vfill~\\

%=====================================================================
\chapter*{Future Development}%
%=====================================================================
\label{ch:future}%
\section*{To be studied$\ldots$}%
\label{ch:future:to-be-studied}%

Following up the current status of the project, the next steps includes:
\begin{enumerate}
  \item Finalization details of Boltzmann Machines, such as the update equations of the state of a unit and the description of the BM algorithm.
  \item Studying Restricted Boltzmann Machines and how its learning algorithm works.
  \item Dummy test to compare results of Hopfield, Boltzmann Machine and Restricted Boltzmann Machine
  \item Identify an interesting problem where the application of RBM is suitable.
\end{enumerate}

\subsection*{Chronogram}%
\label{ch:future:chronogram}

Aiming to conclude this masters dissertation the chronogram below seems to be the most coherent:
\begin{itemize}
  \item Finalization of both Boltzmann and Restricted Boltzmann Machines by the end of \textbf{January 2020}.
  \item Definition of the case to apply RBM by the end of \textbf{January 2020}.
  \item Implementation and testing of RBM algorithm to solve proposed task by the end of \textbf{May 2020}.
  \item Meanwhile keep with the development of the written text in parallel which will take up to about \textbf{June 2020}.
\end{itemize}


%=====================================================================
%\chapter*{Conclusion}
%=====================================================================

%Texto da conclusão.

% ----------------------------------------------------------
% ELEMENTOS POS-TEXTUAIS
% ----------------------------------------------------------
\backmatter%
%=====================================================================
% Referencias via BibTeX
%=====================================================================
\citeoption{abnt-options4.bib}
\bibliography{bibs/bibliografia.bib}
%=====================================================================

%=====================================================================
%\postextualchapter*{Glossário}
%=====================================================================
%\definicao{termo}{significado}
%\definicao{termo}{significado}
%\definicao{termo}{significado}
% ----------------------------------------------------------
% Apêndices (opcionais)
% ----------------------------------------------------------
% ---
% Inicia os apêndices
% ---
\appendix
%=====================================================================
\postextualchapter{Kullback-Leibler Divergence or Relative Entropy}%
%=====================================================================
\label{app:dkl}%
\section{Kullback-Leibler Divergence}%
\label{app:dkl:dkl}%

The Kullback-Leibler divergence or relative entropy is the measure of the difference between two probability distributions. 
The entropy computes how uncertainty events are. 
If a system is quite stable, the entropy will be close to zero and will not vary a lot, while if the system is quite unstable, then the entropy will be very large. 



%\section{Primeira seção}
%
%Texto da primeira seção.
%
%\subsection{Primeira subseção}
%
%Texto da primeira subseção.
%
%\subsubsection{Primeira subsubseção}
%
%Texto da primeira subsubseção.
%=====================================================================
%\postextualchapter{Segundo apêndice}
%=====================================================================
%\section{Primeira seção}
%
%Texto da primeira seção.
%
%\subsection{Primeira subseção}
%
%Texto da primeira subseção.
%
%\subsubsection{Primeira subsubseção}
%
%Texto da primeira subsubseção.
% ----------------------------------------------------------
% Anexos (opcionais)
% ----------------------------------------------------------
% ---
% Inicia os anexos
% ---
\annex%
%=====================================================================
%\postextualchapter{Kullback-Leibler Divergence or Relative Entropy}
%=====================================================================
%\section{Primeira seção}
%
%Texto da primeira seção.
%
%\subsection{Primeira subseção}
%
%Texto da primeira subseção.
%
%\subsubsection{Primeira subsubseção}
%
%Texto da primeira subsubseção.
%=====================================================================
%\postextualchapter{Segundo anexo}
%=====================================================================
%\section{Primeira seção}
%
%Texto da primeira seção.
%
%\subsection{Primeira subseção}
%
%Texto da primeira subseção.
%
%\subsubsection{Primeira subsubseção}
%
%Texto da primeira subsubseção.
%---------------------------------------------------------------------
% INDICE REMISSIVO (relativo ao makeindex)
%---------------------------------------------------------------------
\printindex
%=====================================================================
\end{document}
